\chapter{三个等大块与含零收尾}\label{sec:blocks}

后面的矩阵构造只负责造出一个零和 $p$ 元组。
本讲先证明，一旦做到这件事，确实能够完成整个配置。
为此需要独立解决 $3p$ 个位置。这个维数适合分成三个等大块，
分组矩阵的行列和都等于 $p$，于是能用非常具体的剪切处理。

\section{三个等大块上的剪切}


若三个块的值为 $a,b,c$，一个 $3\times3$ 非负整数矩阵 $K$，
其每行每列的和都为 $p$，就描述一个真实分组：
第 $i$ 个新组从第 $j$ 个旧块取 $K_{ij}$ 个位置。
行和保证每组恰有 $p$ 项，列和保证每个旧位置恰好用一次。
对三个新组分别平均，新块值就是 $K(a,b,c)^T/p$。

下面的 $mJ$ 项在 $a+b+c=0$ 的平面上消失。
留下来的 $I\pm D$ 才是我们想要的剪切；添加 $mJ$ 是为了让计数全部非负。


本讲使用 $\SL_2(\Z)$ 表示行列式为 $1$ 的二阶整数矩阵群。
在两参数平面上，上剪切 $U(k)$ 把 $(x,y)$ 改为 $(x+ky,y)$，
下剪切 $L(k)$ 把它改为 $(x,y+kx)$。
矩阵群的一般同余工具将放到第\ref{sec:arithmetic}--\ref{ch:congruence}讲；
这里的层为 $3$，只需欧几里得算法和四个陪集的直接计算。

\begin{proposition}[3p个位置的独立解法]\label{prop:threep}
若 $p\ge5$ 为素数，则 $3p$ 个位置满足定理\ref{thm:prime}的判据。
\end{proposition}
\begin{proof}
先将位置分为三个 $p$ 组。若其和模 $3$ 全相等，交换不同组中一对不同余值，组和变为 $s+\delta,s-\delta,s$，其中 $\delta\ne0$。这样的跨组数对必存在，否则原配置模 $3$ 全相等。分别平均后得到 $(a^p,b^p,c^p)$，$a+b+c=0$，且模 $3$ 不全相等。

写 $p=3m+\epsilon$，$\epsilon=\pm1$、$m\ge2$。令 $J$ 为三阶全一矩阵，
\[
D=\begin{pmatrix}1&-1&0\\1&-1&0\\-2&2&0\end{pmatrix},\qquad
K_\pm=mJ+\epsilon(I\pm D).
\]
两表非负且行列和为 $p$。每个输出组按一行从三个旧块取数，三次平均实现 $K_\pm/p$。在零和空间上，它们等于 $\epsilon(I\pm D)/p$。由于 $D^2=0$，有
\[
(I+D)(I-D)=I-D^2=I.
\]
实际三步平均仍会带上 $1/p$ 的共同尺度，但两作用在共同尺度意义下互为逆。
因此两个方向的整数次剪切都有正向平均序列。
改用 $(x,y)=(a,a-b)$，它们分别为 $U(1),U(-1)$。块角色交换给
\[
\Sigma=\begin{pmatrix}1&-1\\0&-1\end{pmatrix},\qquad
P=\begin{pmatrix}-2&1\\-3&2\end{pmatrix},\qquad
L(-3)=U(-1)P\Sigma.
\]
因此 $U(\Z),L(3\Z)$ 可实现。


这里 $\Gamma_0(3,\Z)$ 表示行列式为 $1$、左下项被 $3$ 整除的整数矩阵群。
为什么恰好出现 $3$？因为我们正在控制三个等大块，
其模 $3$ 的差分就是剩余的共同余数障碍。
这些矩阵射影生成 $\Gamma_0(3,\Z)$。为验证，取 $S=\left(\begin{smallmatrix}0&-1\\1&0\end{smallmatrix}\right)$、$T=U(1)$，右陪集代表为 $I,ST^{-1},S,ST$，由模 $3$ 底行的四个射影类区分。
陪集表的含义是：从代表 $r$ 出发，右乘一个生成元 $g$，
再乘新陪集代表 $r'$ 的逆，便回到子群。
任意群词可以逐字这样改写为这些回路的乘积，
因此只要表中的回路都由我们已有的矩阵生成，整个子群就由它们生成。
对生成元 $S,T,T^{-1}$，Schreier 回路 $rg(r')^{-1}$ 如下：
\[
\begin{array}{c|ccc}
r&S&T&T^{-1}\\\hline
I&I&T&T^{-1}\\
ST^{-1}&-T&I&L(3)\\
S&-I&I&I\\
ST&T^{-1}&L(-3)&I
\end{array}
\]
欧几里得消元证明 $S,T$ 生成 $\SL_2(\Z)$，上述表便给所需子群；中央符号不影响式\eqref{eq:return}。

本原三块的 $G=\gcd(y,3)$。合法性给 $3\nmid y$，可在 $\alpha x+\beta y=1$ 的裴蜀解中令 $3\mid\alpha$。矩阵
\[
\begin{pmatrix}y&-x\\\alpha&\beta\end{pmatrix}\in\Gamma_0(3,\Z)
\]
将 $(x,y)$ 送到 $(0,1)$，故得到 $(0^p,s^p,(-s)^p)$。
到了相反值配置，令 $h=(p-1)/2$。第一次选 $h$ 份 $s$、$h$ 份 $-s$ 和一个零，
得到 $p$ 个新零；第二次作同样的选择以后，每个符号只剩一份。
最后一组用这两份相反值和 $p-2$ 个零即可。
具体地，连续两次平均 $(s^h,(-s)^h,0)$，最后平均 $(s,-s,0^{p-2})$。
\end{proof}


\section{一个零为什么已经有用}
在 $2p+1$ 个位置中保留一个零，其他位置正好分成两个 $p$ 组。
零和要求两个组的均值互为相反数。因此，上一个证明末尾的三步操作可以直接接上。

\begin{lemma}[含零的2p+1元配置]\label{lem:zero-base}
对奇数 $p$，$2p+1$ 元零和配置若含一个零，则至多五次 $p$ 平均归零。
\end{lemma}
\begin{proof}
保留零，将其余位置分成两个 $p$ 组并平均，得到 $(a^p,(-a)^p,0)$。用命题\ref{prop:threep}末尾所列的三次相反值操作完成。
\end{proof}


\section{按非零支持数选择已经解决的子系统}
一个零和 $p$ 元组平均以后，至少产生 $p$ 个零，所以非零位置最多还有 $n-p$ 个。
我们只在一个包含\emph{全部}非零位置的子系统上工作；外部留下的全是零，
故子系统仍然零和，也不会有缩放尺度不同步的问题。

若一个整数配置含零，坐标差的最大公因子就等于坐标本身的最大公因子。
本原化以后它的 $G$ 必为 $1$。因此这个子系统可以直接调用已证维数的充分性。
选择子系统时只需让位置数严格大于非零支持数，确保其中也有一个零。

\begin{lemma}[从一个零和组到全部归零]\label{lem:zero-tail}
若 $p\ge5$，$n\ge2p+1$ 的零和配置含一个零和 $p$ 元组，则可归零。对任意奇数 $p$，此结论在 $2p+1\le n\le3p$ 同样成立。
\end{lemma}
\begin{proof}
先平均该组，产生 $p$ 个零。若某个已解决的维数 $m\le n$ 满足非零支持数小于 $m$，就选一个 $m$ 元子系统包含全部非零位置及一个零。它自身零和，本原化后的差分公因子为 $1$，可在该子系统归零。
对 $n\le3p$，使用引理\ref{lem:zero-base}中的 $m=2p+1$。对 $3p<n<4p$，使用命题\ref{prop:threep}。对 $4p\le n<5p$，使用引理\ref{lem:integer-range}在 $m=4p$ 的结论，因为 $d_p(4p)=1$。对 $n\ge5p$，同一引理的最终上界已适用。
\end{proof}

